\chapter{Hiện thực và đánh giá}
\label{chap:implementation}

\section{Các bước triển khai}

Đây là nơi phù hợp để mô tả nhóm đã hiện thực đề tài như thế nào: khởi tạo môi trường,
chia module, cấu hình hạ tầng, và kết nối giữa các thành phần.

\begin{lstlisting}[language=Python, caption={Ví dụ đoạn mã minh họa cách chèn code.}, label={lst:demo}]
from fastapi import FastAPI

app = FastAPI()

@app.get("/health")
def healthcheck():
    return {"status": "ok"}
\end{lstlisting}

Listing \ref{lst:demo} cho thấy template đã sẵn sàng cho việc chèn code có caption và label.

\section{Kết quả và đánh giá}

Cần nêu rõ nhóm đánh giá bằng cách nào: testcase, user testing, benchmark, hay so sánh với yêu cầu ban đầu.
Nếu có biểu đồ, ảnh chụp giao diện, hoặc bảng số liệu, chèn vào đây để tăng độ thuyết phục.

\begin{table}[H]
  \centering
  \begin{tabular}{|>{\bfseries}l|c|c|}
    \hline
    Tiêu chí & Mục tiêu & Kết quả \\
    \hline
    Tỷ lệ testcase đạt & 90\% & 96\% \\
    \hline
    Thời gian phản hồi trung bình & < 500 ms & 320 ms \\
    \hline
    Tỷ lệ hoàn thành backlog & 100\% & 92\% \\
    \hline
  \end{tabular}
  \caption{Mẫu bảng tổng hợp kết quả đánh giá.}
  \label{tab:evaluation}
\end{table}

Bảng \ref{tab:evaluation} là chỗ đặt sẵn để thay bằng số liệu thật của đồ án.
Nếu cần thêm phần trích dẫn tài liệu tham khảo, có thể dùng \cite{mittelbach2004latexcompanion}.
